% ==============================================================================
% CONCLUSIONES
% ==============================================================================
% Solo componente de ciberseguridad (decisión del estudiante, 2026-10-09):
% no se incluyen conclusiones sobre entrenamiento o selección de modelos,
% que corresponden al componente de IA.
% Cada conclusión se deriva de una cifra de k-resultados.tex y de su
% interpretación en l-discusion.tex.
\ifdefined\modoensamble
\else
	\documentclass[11pt, letterpaper]{report}
\usepackage[spanish, es-nodecimaldot, es-noquoting]{babel}
\selectlanguage{spanish}
\usepackage[utf8]{inputenc}
\usepackage[T1]{fontenc}

% ===============================================================================
% El estudiante debe llenar sus datos en esta sección para que la plantilla los 
% auto-importe y genere automáticamente las páginas de portada y de firmas 
% autorizadas.
% ===============================================================================
% Datos del estudiante:
% -------------------------------------------------------------------------------
% Nombre completo
\def \nombreestudiante {Sebastián Huertas Gómez}
% Carné
\def \uvgcarne {22295}
% Facultad
\def \uvgfacultad {Ingeniería}
% Carrera
\def \uvgcarrera {Ingeniería en Ciencia de la Computación y Tecnologías de la Información}

% Datos del trabajo:
% -------------------------------------------------------------------------------
% Título completo (tomado literalmente de protocolo_sebastian_v1 (3).pdf, sección 1)
\def \titulotesis {logSguarDian: Desarrollo de una librería para la detección de amenazas de seguridad en Endpoints de aplicaciones}
% Mes y año de entrega (confirmado por el estudiante: noviembre 2026)
\def \mesentrega {noviembre}
\def \anoentrega {2026}
% Asesor (según protocolo_sebastian_v1 (3).pdf, portada)
\def \nombreasesor {Ing. Erick Francisco Marroquín Rodríguez}
% El catedrático del curso (Ing. Bacilio Alexander Bolaños Lima) figura como
% primer examinador en la hoja de firmas, definido más abajo.

% Datos del tribunal examinador:
% -------------------------------------------------------------------------------
% Nombre del primer examinador (segunda línea de la hoja de firmas)
\def \nombreprimerex {Ing. Ma. Bacilio Alexander Bolaños Lima}

% Capítulos pre-definidos
% -------------------------------------------------------------------------------
% Comentar las líneas de las secciones que desean omitirse, por defecto se
% se incluyen todas.
% Se descomentan conforme se redacta cada capítulo (todavía sin archivo de
% contenido: prefacio, alcance, anexos, glosario).
\def \CAPprefacio {Prefacio}
\def \CAPantecedentes {Antecedentes}
\def \CAPalcance {Alcance}
\def \CAPmetodologia {Metodología}
\def \CAPresultados {Resultados}
\def \CAPdiscusion {Análisis de resultados}
\def \CAPanexos {Anexos}
\def \CAPglosario {Glosario}

% Formato y estilo de la plantilla
% -------------------------------------------------------------------------------
% Modo impresión: Puede des-comentar la siguiente línea para generar un documento pdf sin la portada, para cuando se desee imprimir el documento para encuadernación
% NOTA: la plantilla original traía esta línea activa por defecto, lo cual
% suprime la portada verde (main.tex hace \let\CAPportada\undefined cuando
% \printver está definido). Se comenta para que la versión digital, con
% portada, sea la que se genera por defecto.
%\def \printver {Versión del documento para impresión}

% Portada: Puede cambiarse la imagen en la portada al cambiar el nombre del 
% archivo siguiente. NOTA: debe tener la suficiente resolución para cubrir el área
% designada
\def \imagenportada {plantilla/portadacit.jpg}

% Referencias: Puede des-comentar la siguiente línea para utilizar el formato de referencias APA
%\def \usarAPA {Usar formato APA}

% Párrafo: Puede comentar la siguiente línea si desea emplear un formato de 
% párrafo distinto al establecido por defecto
\def \parpordefecto {Formato de párrafo por defecto}

% Capítulos y secciones: Puede des-comentar la siguiente línea para establecer el
% formato de los capítulos y secciones bajo el estándar original de UVG para
% trabajos de graduación. Este incluye: capítulos con numeración romana, secciones
% con letras mayúsculas, sub-secciones con números y sub-sub-secciones con letras
% minúsculas
%\def \capsecuvg {Formato UVG para capítulos y secciones}
% ================================================================================
% En este archivo se colocan opciones adicionales para modificar el formato de la
% plantilla, para emplearse en otros tipos de documentos que no sean trabajos de
% graduación. Si usted está trabajando su tesis, NO modifique este archivo
% ================================================================================
% Capítulos pre-definidos
% --------------------------------------------------------------------------------
% Comentar las líneas de las secciones que desean omitirse, por defecto se 
% se incluyen todas.
\def \CAPportada {Portada}
\def \CAPcaratula {Caratula}
\def \CAPfirmas {Hoja de firmas}
\def \CAPindice {Índice general}
\def \CAPfiguras {Listado de figuras}
\def \CAPcuadros {Listado de cuadros}
\def \CAPresumen {Resumen}
\def \CAPabstract {Resumen}
\def \CAPintroduccion {Introducción}
\def \CAPobjetivos {Objetivos}
\def \CAPjustificacion {Justificación}
\def \CAPmarcoteorico {Marco teórico}
\def \CAPconclusiones {Conclusiones}
\def \CAPrecomendaciones {Recomendaciones}
\def \CAPbibliografia {Bibliografía}

\usepackage{xcolor}
\usepackage{amsfonts}
\usepackage{amsmath}
\usepackage{amssymb}
\usepackage{amsthm}
\usepackage{mathtools}
\usepackage{graphicx}
\usepackage{xfrac}
\usepackage{float}
\usepackage[hypertexnames=false]{hyperref}
\usepackage[font=small]{caption}
\usepackage{subcaption}
% Cuadros: array (columnas L{} alineadas a la izquierda) y longtable (cuadros
% que ocupan más de una página). Se cargan aquí, no en
% 2-paquetes_y_comandos_usuario.tex, para que los capítulos de Metodología y
% Resultados no dependan de ese archivo.
\usepackage{array}
\usepackage{longtable}
\newcolumntype{L}[1]{>{\raggedright\arraybackslash}p{#1}}
\renewcommand{\arraystretch}{1.3}
\setlength{\LTpre}{0.3cm}
\setlength{\LTpost}{0.1cm}
% Margen de estiramiento para que los párrafos con palabras largas no
% generen líneas que sobresalgan del margen.
\setlength{\emergencystretch}{3em}
\usepackage{xpatch}
\usepackage{emptypage}
% htt: permite dividir con guion el texto en \texttt (nombres de archivo,
% clases y campos), que de otro modo sobresale del margen derecho.
\usepackage[htt]{hyphenat}
\usepackage{fancyhdr}
\usepackage[backend=biber, style=ieee]{biblatex}
\ifdefined\usarAPA
    \usepackage[backend=biber, style=apa]{biblatex}
\fi
\addbibresource{m-bibliografia.bib}

\usepackage[percent]{overpic}
\usepackage{chngcntr}

\ifdefined\CAPglosario
	\usepackage[numberedsection]{glossaries}
	\makeglossaries
    \newglossaryentry{rasp}
{
    name={RASP},
    description={Autoprotección de aplicaciones en tiempo de ejecución (\textit{Runtime Application Self-Protection}): mecanismo de seguridad que se ejecuta dentro de la propia aplicación y analiza sus peticiones para detectar y bloquear ataques.}
}

\newglossaryentry{waf}
{
    name={WAF},
    description={Cortafuegos de aplicaciones web (\textit{Web Application Firewall}): dispositivo o servicio perimetral que filtra el tráfico HTTP antes de que llegue a la aplicación.}
}

\newglossaryentry{modsecurity}
{
    name={ModSecurity},
    description={Motor de WAF de código abierto que se usa con reglas externas; en este trabajo se configuró con OWASP CRS.}
}

\newglossaryentry{crs}
{
    name={OWASP CRS},
    description={Conjunto de reglas genéricas de OWASP (\textit{Core Rule Set}) para ModSecurity, con niveles de paranoia que controlan qué tan estrictas son.}
}

\newglossaryentry{owasp}
{
    name={OWASP Top 10},
    description={Lista de las diez categorías de riesgo más críticas para aplicaciones web, publicada por OWASP; este trabajo usa la edición 2021.}
}

\newglossaryentry{mitre}
{
    name={MITRE ATT\&CK},
    sort={MITRE ATT&CK},
    description={Base de conocimiento de tácticas y técnicas de atacantes, usada aquí para asociar cada vector de ataque a una técnica.}
}

\newglossaryentry{cwe}
{
    name={CWE},
    description={Enumeración común de debilidades de software (\textit{Common Weakness Enumeration}): catálogo de tipos de defectos que originan vulnerabilidades.}
}

\newglossaryentry{sqli}
{
    name={SQL Injection},
    description={Ataque que inserta sentencias SQL en una entrada de la aplicación para alterar las consultas que ejecuta la base de datos.}
}

\newglossaryentry{xss}
{
    name={XSS},
    description={\textit{Cross-Site Scripting}: ataque que inyecta código ejecutable en una página que otros usuarios verán.}
}

\newglossaryentry{pathtraversal}
{
    name={Path Traversal},
    description={Ataque que manipula rutas de archivo (por ejemplo con \texttt{../}) para leer archivos fuera del directorio previsto; incluye la inclusión local de archivos (LFI).}
}

\newglossaryentry{cmdi}
{
    name={Command Injection},
    description={Ataque que logra ejecutar comandos del sistema operativo a través de una entrada de la aplicación.}
}

\newglossaryentry{middleware}
{
    name={\textit{middleware}},
    sort={middleware},
    description={Función que se ejecuta entre la recepción de una petición y la lógica de negocio de la aplicación; en Express se registra con \texttt{app.use}.}
}

\newglossaryentry{rf}
{
    name={Random Forest},
    description={Algoritmo supervisado que combina muchos árboles de decisión; en la librería es el único modelo con autoridad para bloquear una petición.}
}

\newglossaryentry{iforest}
{
    name={Isolation Forest},
    description={Algoritmo no supervisado de detección de anomalías que aísla las observaciones atípicas con particiones aleatorias; en la librería solo registra anomalías y nunca bloquea.}
}

\newglossaryentry{onnx}
{
    name={ONNX},
    description={Formato abierto para serializar modelos de aprendizaje automático, que permite ejecutarlos fuera de Python.}
}

\newglossaryentry{workerthreads}
{
    name={\texttt{worker\_threads}},
    sort={worker_threads},
    description={Módulo de Node.js que ejecuta código en hilos independientes del hilo principal; la inferencia de los modelos se ejecuta ahí.}
}

\newglossaryentry{eventloop}
{
    name={\textit{Event Loop}},
    sort={Event Loop},
    description={Bucle de eventos de Node.js que atiende las peticiones en un único hilo; bloquearlo degrada a toda la aplicación.}
}

\newglossaryentry{failopen}
{
    name={\textit{fail-open}},
    sort={fail-open},
    description={Política por la cual, ante un error o un tiempo de espera del módulo de detección, la petición se deja pasar en lugar de bloquearse.}
}

\newglossaryentry{monitor}
{
    name={Modo \texttt{monitor}},
    sort={Modo monitor},
    description={Modo de operación en el que la librería registra las peticiones que bloquearía, pero no las bloquea.}
}

\newglossaryentry{block}
{
    name={Modo \texttt{block}},
    sort={Modo block},
    description={Modo de operación en el que la librería responde HTTP 403 a las peticiones que clasifica como ataque.}
}

\newglossaryentry{pesp}
{
    name={p95},
    description={Percentil 95 del tiempo de respuesta: el valor por debajo del cual queda el 95~\% de las peticiones.}
}

\newglossaryentry{deltap95}
{
    name={$\Delta p95$},
    sort={Delta p95},
    description={Diferencia entre el p95 con el middleware activo y el p95 de la línea base sin él; en este trabajo se mide en milisegundos.}
}

\newglossaryentry{funo}
{
    name={F1},
    description={Media armónica de la precisión y el recall; es la métrica principal de clasificación del trabajo.}
}

\newglossaryentry{auc}
{
    name={AUC-ROC},
    description={Área bajo la curva ROC; mide qué tan bien separa un modelo dos clases sin depender de un umbral.}
}

\newglossaryentry{wilson}
{
    name={Intervalo de Wilson},
    description={Intervalo de confianza para una proporción, adecuado para muestras pequeñas o proporciones cercanas a 0 o 1.}
}

\newglossaryentry{casidup}
{
    name={Casi-duplicado},
    description={Registro muy parecido a otro pero no idéntico; si queda repartido entre entrenamiento y prueba, puede inflar las métricas.}
}

\newglossaryentry{minhash}
{
    name={MinHash},
    description={Técnica para estimar rápidamente la similitud entre conjuntos, usada para detectar casi-duplicados.}
}

\newglossaryentry{fuga}
{
    name={Fuga de datos},
    description={Presencia en el conjunto de evaluación de información disponible durante el entrenamiento, lo que sobreestima el desempeño.}
}

\newglossaryentry{corpuslimpio}
{
    name={Corpus limpio},
    description={Conjunto de payloads de evaluación filtrado para excluir los que coinciden con el entrenamiento; en este trabajo es el de la Ronda~5.}
}

\newglossaryentry{useragent}
{
    name={\texttt{User-Agent}},
    sort={User-Agent},
    description={Cabecera HTTP con la que el cliente se identifica (navegador, cURL, entre otros); su representación en el corpus sesgó al modelo.}
}

\newglossaryentry{falsopositivo}
{
    name={Falso positivo},
    description={Petición legítima que el sistema clasifica o bloquea como ataque.}
}

\newglossaryentry{parmodelos}
{
    name={Par de modelos},
    description={Combinación de un archivo \texttt{rf.onnx} y uno \texttt{if.onnx} identificada por su suma SHA-256 (pares A, B y C).}
}

\newglossaryentry{tarball}
{
    name={\textit{tarball}},
    sort={tarball},
    description={Archivo comprimido \texttt{.tgz} de un paquete de npm, publicado o construido localmente.}
}

\newglossaryentry{kanban}
{
    name={Kanban},
    description={Metodología de gestión basada en un tablero de columnas que representan el estado de las tareas.}
}

\newglossaryentry{artillery}
{
    name={Artillery},
    description={Herramienta de pruebas de carga usada para medir la latencia con y sin el middleware.}
}

\newglossaryentry{sqlmap}
{
    name={sqlmap},
    description={Herramienta de código abierto que automatiza la detección y la explotación de inyecciones SQL; generó los payloads de SQL Injection del corpus limpio.}
}

\newglossaryentry{zap}
{
    name={OWASP ZAP},
    description={Herramienta de OWASP para pruebas de seguridad de aplicaciones web; se usó para confirmar el payload de XSS.}
}

\newglossaryentry{mcp}
{
    name={MCP},
    description={Protocolo de Contexto de Modelo (\textit{Model Context Protocol}): estándar para que un asistente de inteligencia artificial use herramientas externas.}
}

\fi

\usepackage[top=1in, left=1.5in, right=1in, bottom=1in]{geometry}
\usepackage[Sonny]{fncychap}

\graphicspath{ {figuras/} }
\definecolor{uvg-green}{RGB}{17,71,52}
\newcommand{\defaultparformat}[1]{
	{\setlength{\parskip}{2ex}
     \input{#1}}
}
\ifdefined\capsecuvg
	\renewcommand\thechapter{\Roman{chapter}}
    \renewcommand\thesection{\Alph{section}}
	\renewcommand\thesubsection{\arabic{subsection}}
    \renewcommand\thesubsubsection{\alph{subsubection}}
\fi
\counterwithout{figure}{chapter}
\counterwithout{table}{chapter}
\counterwithout{equation}{chapter}

\providecommand{\blankpage}{
\newpage
\thispagestyle{empty}
\mbox{}
\newpage
}

% ==============================================================================
% Paquetes y comandos adicionales del usuario
% ==============================================================================
% TikZ: requerido por las figuras del capítulo de Marco Teórico (pipeline de
% middleware, comparación RASP vs. control perimetral).
\usepackage{tikz}
\usetikzlibrary{arrows.meta, positioning, shapes.geometric, fit, backgrounds}


\pagestyle{headings}
\begin{document}

\ifdefined\parpordefecto
	\setlength{\parskip}{2ex}
\fi

\fi

\chapter{Conclusiones}
\label{cap:conclusiones}

Las conclusiones siguen el orden de los objetivos: primero el objetivo general, después los
objetivos específicos y, al final, la hipótesis y la contribución al problema. Cada una indica el
objetivo al que responde, el grado de cumplimiento y el alcance de la evidencia que la sostiene.

\begin{enumerate}
    \item \textbf{Objetivo general: cumplido de forma parcial.}\label{conc:og} Se diseñó, implementó y publicó en
    el registro público de npm la librería \texttt{logsguardian} (versión 0.1.0, 30 de agosto de
    2026, y versión 0.1.1, 11 de octubre de 2026), un middleware RASP para Express que integra un
    extractor de 76 características y un modelo híbrido (Random Forest e Isolation Forest)
    serializado en ONNX y ejecutado en hilos de trabajo, con una diferencia máxima de
    $2.42\times10^{-7}$ entre las predicciones de Python y de ONNX. La librería intercepta las
    peticiones, clasifica las cuatro categorías de ataque y bloquea con el Random Forest. No se
    cumplen los criterios de memoria, CPU ni carga de volumen (conclusión~\ref{conc:rendimiento}), y
    el aporte del Isolation Forest a la detección de variantes desconocidas solo tiene evidencia
    indirecta (conclusión~\ref{conc:variantes}). La versión 0.1.0 contiene la generación de modelo
    \texttt{rf\_v10}/\texttt{if\_v9}; la versión 0.1.1, vigente, contiene \texttt{rf\_v11}/\texttt{if\_v10},
    que es la evaluada sobre el conjunto de prueba, y la corrección de un defecto del middleware.

    \item \textbf{Objetivo específico 1: cumplido.} Los cuatro vectores de ataque intra-request se
    mapearon a OWASP Top 10 2021 (SQL Injection, XSS y Command Injection en A03:2021; Path
    Traversal/LFI en A01:2021) y a MITRE ATT\&CK (T1190 y T1189), y ese mapeo estructuró las cinco
    clases canónicas del dataset y las categorías sobre las que se calculan las métricas del
    objetivo específico~3.

    \item \textbf{Objetivo específico 2: cumplido.} Se construyó un corpus de 383{,}202 peticiones
    HTTP etiquetadas, 3.8 veces el mínimo de 100{,}000, a partir de siete fuentes públicas
    auditadas, dos listas públicas de payloads y conjuntos propios, con etiquetado determinista y
    una partición verificada y bloqueada por hash. El desbalance de clases (59.3~\% SQLi frente a
    2.6~\% Command Injection) no impidió que la clase minoritaria alcanzara un F1 de 0.9239.

    \item \textbf{Objetivo específico 3, clasificación: cumplido.}\label{conc:clasificacion} Sobre el conjunto de prueba
    bloqueado, el sistema obtuvo un F1 de entre 0.9239 (Command Injection) y 0.9956 (SQL Injection)
    en las cuatro categorías, frente al criterio de 0.80 en al menos tres, y el Isolation Forest
    detectó el 91.25~\% de los ataques con 5.46~\% de falsos positivos, frente a los criterios de
    al menos 50~\% y como máximo 10~\%. Estas cifras miden la separación entre clases dentro del
    corpus y no la generalización a tráfico externo: el conjunto de prueba se leyó tres veces en la
    generación vigente y el corpus tiene saturación de casi-duplicados, en particular en SQL
    Injection (unos 395 pares casi duplicados por fila), un efecto que este trabajo no cuantificó.

    \item \textbf{Objetivo específico 3, cobertura sobre payloads no vistos: cumplido, con reservas.}\label{conc:cobertura}
    Sobre 600 payloads filtrados contra el entrenamiento, la cobertura fue, con el paquete 0.1.0
    (tarball de la Ronda~5), de 100~\% en SQL Injection, 85.0~\% en Path Traversal y 91.5~\% en
    Command Injection; con el tarball del modelo vigente, de 99.5~\%, 95.0~\% y 96.0~\%; y con el
    paquete 0.1.1 y el ejecutor de pruebas corregido, de 99.5~\%, 100.0~\% (200 de 200; intervalo de
    Wilson al 95~\% de 98.1 a 100.0~\%) y 92.5~\%. Los fallos de Path Traversal de la primera
    medición tenían dos causas independientes, un defecto del middleware y otro del ejecutor, ambos
    corregidos; la cobertura de Command Injection bajó de 96.0~\% a 92.5~\% porque la cifra anterior
    incluía 10 payloads que solo se bloqueaban en su forma corrompida. Sobre un corpus propio de 200
    payloads de XSS, el paquete 0.1.1 bloqueó 200 de 200 (intervalo de Wilson al 95~\% de 98.1 a
    100.0~\%; 64 de 64 en el campo vulnerable), de modo que el criterio del 80~\% se cumple para
    XSS en cobertura. Las reservas son dos: la medición con el paquete 0.1.1 de Path Traversal y de
    Command Injection se hizo sobre el corpus que reveló los defectos, por lo que no es
    independiente; y el corpus de XSS se construyó después de conocer los resultados de la Ronda~5,
    136 de sus 200 payloads apuntan a campos sin vulnerabilidad confirmada, y la explotabilidad
    confirmada sigue descansando en un único caso de OWASP ZAP (n=1).

    \item \textbf{Objetivo general, rendimiento del servidor anfitrión: no cumplido en su
    totalidad.}\label{conc:rendimiento} El extractor cumple el criterio en aislamiento: el p95 con
    \texttt{worker\_threads} es de 0.0273~ms, frente al límite de 5~ms. El sobrecosto del middleware
    completo es de 12.08~ms en carga normal y de 12.85~ms en carga de ataque (límite de 15~ms), y la
    extracción no lo explica. El extractor inicial de 72 características y el paso a
    \texttt{worker\_threads} los implementó el autor de este trabajo; el análisis y la justificación
    de cada característica, el contrato de entrada, las cuatro características adicionales y las
    correcciones del extractor corresponden al componente de Inteligencia Artificial. El criterio
    de latencia no se cumple en el escenario de volumen, donde el $\Delta p95$ fue de 6{,}495.46~ms
    y el 7.00~\% de las peticiones agotó el tiempo de espera con el middleware activo. Los
    indicadores de recursos del protocolo tampoco se cumplen: el middleware agregó entre 331.9 y
    375.9~MiB de memoria según el escenario, más del doble del límite de 150~MB (143.1~MiB), y
    +529.6~\% a +614.4~\% de CPU, frente a un límite de 15~\%. El incremento de memoria es casi
    constante entre escenarios, pero su reparto por componente no se midió.

    \item \textbf{Objetivo general, variantes desconocidas e Isolation Forest.}\label{conc:variantes}
    La evidencia directa de que la librería bloquea variantes desconocidas de los vectores cubiertos
    es la cobertura del Random Forest descrita en la conclusión~\ref{conc:cobertura} (4 de los 200 payloads de SQL
    Injection coinciden por igualdad normalizada con alguna partición; ninguno de Path Traversal ni
    de Command Injection). La contribución del Isolation Forest solo tiene evidencia indirecta: en
    el conjunto de validación del componente de Inteligencia Artificial \cite{valenzuela2026ia} marca 19 de los 39 ataques
    que el Random Forest deja pasar, pero bloquear sus marcas agregaría 866 benignos bloqueados de
    más, y por eso solo registra. Esas cifras son optimistas, porque el umbral del Isolation Forest
    se eligió sobre el mismo conjunto de validación, y no se evaluaron sobre el conjunto de prueba.

    \item \textbf{Objetivo general, defensa en profundidad.}\label{conc:waf} En la Ronda~4, ModSecurity con OWASP CRS
    bloqueó el 93.2~\% de 590 payloads, la librería el 98.8~\% y ambas capas el 99.5~\%; la librería
    bloqueó 37 de los 40 ataques que el WAF dejó pasar, ninguno de ellos de Command Injection, la
    categoría con fuga de datos confirmada en ese corpus. En la Ronda~5, sobre un corpus sin esa fuga y ampliado
    con 200 payloads de XSS, medido el 25 de octubre de 2026 con las cuatro configuraciones sobre
    el mismo árbol de la aplicación y el ejecutor corregido, la librería sola bloqueó 784 de 800
    payloads, el WAF solo 749 y ambas capas 797. El WAF bloqueó el 82.0~\% de Command Injection
    (intervalo de Wilson al 95~\%: 76.1--86.7~\%) frente al 92.5~\% de la librería, y el 68.0~\%
    reportado antes se debió al ejecutor anterior y no a la aplicación. De los 51 ataques que el WAF
    no bloqueó, ambas capas bloquearon 48 y la librería sola 44. La complementariedad es asimétrica:
    la combinación agrega 13 payloads a la librería y 48 al WAF, en Path Traversal y en XSS la
    librería sola iguala a la combinación, y solo en SQL Injection y en Command Injection cada capa
    bloquea payloads que la otra no. La evidencia respalda el uso de la librería como capa
    complementaria a un WAF y no como su sustituto, con la limitación de que el corpus se construyó
    a partir de familias de payload representadas en el entrenamiento de la librería y de que no se
    midieron otros niveles de paranoia del WAF, por lo que la comparación no es neutral
    (Sección~\ref{sec:met-validez}).

    \item \textbf{Objetivo general, falsos positivos y generalización.}\label{conc:fp} Las métricas del corpus no
    predijeron el comportamiento con tráfico legítimo real, y el corpus benigno es la debilidad
    principal. Con un \texttt{User-Agent} de navegador, la librería en modo \texttt{block} bloqueó el
    92.5~\% del tráfico legítimo en dvws-node (37 de 40; el modelo exacto de esa medición no consta),
    con muestras preliminares en DVNA (5 de 5) y de una sola ruta en node-api-goat (50 de 50). En
    OWASP Juice Shop, en modo \texttt{monitor}, el modelo asignó veredicto de bloqueo al 97.1~\% de
    las peticiones legítimas (303 de 312). En la aplicación de referencia, con otro modelo, ninguna
    petición legítima de Artillery recibió un 403, y con cadenas de \texttt{User-Agent} de
    navegadores reales se bloquearon 13 de 15 combinaciones, frente al 0.15~\% del conjunto de
    prueba. La exclusión por fuente del componente de Inteligencia Artificial \cite{valenzuela2026ia} apunta en el mismo
    sentido: al excluir \texttt{modsec\_learn}, \texttt{payloads\_csv} o \texttt{xss\_dataset} del
    entrenamiento, el bloqueo de filas legítimas de esas fuentes sube de 0.1~\%, 0.4~\% y 0.2~\% a
    76.0~\%, 27.0~\% y 94.5~\%. El sesgo de representación de la cabecera \texttt{User-Agent} es la
    explicación principal, pero no la única (dvws-node bloqueó el 75~\% con el \texttt{User-Agent}
    por defecto de Node.js y Juice Shop dio 39.4~\% con cURL) y no se verificó con un corpus legítimo
    alternativo. En su estado actual, la librería no es apta para instalarse en modo de bloqueo sin
    calibración previa con el tráfico de la aplicación.

    \item \textbf{Trazabilidad y validez de las métricas.}\label{conc:trazabilidad} La auditoría de fuga entre los corpus de
    evaluación y de entrenamiento redujo la cobertura atribuible de Command Injection de 100~\% a
    91.5~\%. El filtro de exclusión, el bloqueo del conjunto de prueba con verificación de hash y el
    manifiesto de sumas de verificación fueron condiciones necesarias para que las cifras reflejaran
    capacidad de detección y no memorización. Esos controles no evitaron que el conjunto de prueba
    se leyera tres veces en la generación vigente, que el origen de los datos restaurados quedara
    sin registrar, ni que, hasta la publicación de la versión 0.1.1, el paquete publicado y las
    métricas de prueba correspondieran a generaciones de modelo distintas, lo que limita la
    trazabilidad de las métricas.

    \item \textbf{Hipótesis.} El trabajo no formuló una hipótesis contrastable; el objetivo general
    cumple ese papel como afirmación de factibilidad: que un middleware de aprendizaje automático
    distribuido como paquete npm puede detectar y clasificar ataques web sin comprometer el
    rendimiento del servidor. La evidencia la respalda en la detección y clasificación y en la
    latencia bajo carga normal y de ataque, y es insuficiente o contraria en memoria, CPU, carga de
    volumen y comportamiento con tráfico legítimo real sin calibración, por lo que queda respaldada
    de forma parcial.

    \item \textbf{Contribución al problema.} El problema de partida es la falta de protección
    accesible a nivel de aplicación en organizaciones sin equipo de seguridad dedicado. El trabajo
    aporta una librería de instalación directa desde npm, sin infraestructura externa ni licencias
    de pago, con un extractor de 76 características y un modelo híbrido de procedencia documentada,
    y establece con evidencia dónde la solución funciona (clasificación de las cuatro categorías,
    cobertura de variantes y latencia bajo carga normal y de ataque) y dónde no (recursos, volumen y
    falsos positivos en aplicaciones de terceros). No demuestra que la herramienta sea adecuada para
    producción, y la validación con una startup quedó fuera de la evidencia del proyecto.
\end{enumerate}

\ifdefined\modoensamble
\else
	\section*{Referencias}
	\printbibliography[heading=none]
	\end{document}
\fi
